%===================================================== SECTION: EMPIRICAL SUITE
\section{Empirical Verification Suite}\label{sec:empirical}

Every claim in this report is backed by an archived run log under \texttt{final-log/}.
Twelve runs span five real operating systems, both defense platforms, a 100-port scale
test, and two SSH post-exploitation modes. Rather than summarise them, this section walks
the decisive runs and shows the \emph{actual tool output} — the real proof of the work.
Each excerpt is verbatim from the cited log.

\renewcommand{\arraystretch}{1.2}
\noindent\begin{tabularx}{\textwidth}{@{}L{4.4cm} L{2.3cm} X@{}}
\toprule
\hdrrow \thh{Log} & \thh{Target} & \thh{Demonstrates}\\
\midrule
{\small\texttt{full-chain.log}}        & Win 11 & Full 5-phase chain + live defense severing (\S\ref{sub:fullchain})\\
{\small\texttt{cross-platform-exploit.log}} & Win + Kali & OS-aware exploitation on two OSes (\S\ref{sub:crossplat})\\
{\small\texttt{kali.log}}              & Kali VM & Zero-open-port \& ICMP fingerprinting (\S\ref{sub:closed})\\
{\small\texttt{scan\_100\_ports.log}}  & Win 10/11 & 100-port scale test (\S\ref{sub:scale})\\
{\small\texttt{defense\_win\_test.log}}& Win 11 & Four-vector defense, \texttt{netsh} (\S\ref{sub:defwin})\\
{\small\texttt{defense\_kali\_test.log}}& Kali VM & Four-vector defense, \texttt{iptables} (\S\ref{sub:defkali})\\
{\small\texttt{ssh\_quiet/attack.log}} & Win 11 & SSH brute-force + post-exploitation (\S\ref{sub:ssh})\\
{\small\texttt{ipad.log}, \texttt{phn.log}} & iPad, Android & Darwin \& Linux fingerprints (\S\ref{sub:mobile})\\
\bottomrule
\end{tabularx}
\renewcommand{\arraystretch}{1.0}

%--------------------------------------------------------------------
\subsection{Flagship end-to-end run}\label{sub:fullchain}
\texttt{full-chain.log} drives the complete toolkit against the Windows~11 victim
(\texttt{192.168.0.116}). Phase~1 classifies each port from its raw reply:
\begin{logbox}
   21  OPEN                           TTL=128 WIN=65535
   22  NO RESPONSE / POSSIBLY FILTERED
   80  OPEN                           TTL=128 WIN=65535
   81  NO RESPONSE / POSSIBLY FILTERED
\end{logbox}

Phase~3 then scores the open-port SYN/ACK. Note the honest confidence (62\%) and the full
evidence trail — TTL, window, DF, the IP-ID sequence, and the exact TCP option order:
\begin{logbox}
  Candidate  : Open-Port TCP SYN/ACK
  Likely OS  : Windows          Confidence : 62%
  Version    : Windows 10 / 11 / Server 2016+
  TTL        : 128 (initial 128, ~0 hop(s))    Window : 65535    DF bit : set
  IP ID Seq  : [27057, 27059, 27062] (deltas: [2, 3])   IP ID Class: STRICT_INCREMENTAL
  Options    : ('MSS',1460),('NOP',),('WSCALE',8),('SACK_PERMITTED',),('TIMESTAMP',...)
\end{logbox}

Phase~5's chained FTP-write$\to$web-shell$\to$RCE exploit confirms remote code execution and
cleans up after itself:
\begin{logbox}
  Status  : SUCCESS
  Detail  : RCE confirmed via uploads/shell.py (command: whoami).
    OK login  OK write_check  OK upload  OK http_access  OK rce  OK cleanup (Deleted shell.py)
  RCE output:  Host OS : Windows 11   Executed : whoami
    homed01\turjo
\end{logbox}

The FTP bounce, by contrast, is correctly \hl{detected as mitigated} rather than failing
silently — the modern Windows stack aborts the third-party pivot:
\begin{logbox}
  Status  : FAILED
  Detail  : PORT command rejected or connection aborted
            (500 [WinError 10053] ...). FTP bounce is mitigated by host firewall.
    192.168.0.116:80 -> blocked
\end{logbox}

%--------------------------------------------------------------------
\subsection{OS-aware exploitation across two operating systems}\label{sub:crossplat}
\texttt{cross-platform-exploit.log} runs the \emph{identical} binary against Windows and
Kali back-to-back. The engine fingerprints each host and Phase~4 swaps the payloads
automatically — the clearest proof that reconnaissance drives weaponisation. The two hosts
fingerprint distinctly (Windows 62\% strict-incremental IP-ID; Linux 66\% all-zero IP-ID),
and the same code path produces OS-correct RCE on each:
\begin{logbox}
[WIN  192.168.0.116]  fingerprint -> Windows (62%)   IP ID Class: STRICT_INCREMENTAL
   payload : windows\win.ini + whoami
   RCE     : homed01\turjo
[KALI 192.168.56.101] fingerprint -> Linux (66%)     IP ID Class: ALL_ZERO
   payload : /etc/passwd + id
   RCE     : uid=0(root) gid=0(root) groups=0(root)
\end{logbox}

%--------------------------------------------------------------------
\subsection{Zero-open-port and ICMP fingerprinting}\label{sub:closed}
\texttt{kali.log} proves the engine still works when \emph{nothing is open}. With all three
probed ports closed, the closed-port RST path alone attributes the OS at 76\% — a fallback
the Design Report did not have:
\begin{logbox}
   22  CLOSED   TTL= 64 WIN=0     80  CLOSED   TTL= 64 WIN=0     81  CLOSED   TTL= 64 WIN=0

  Candidate  : Closed-Port TCP RST      Likely OS : Linux      Confidence : 76%
  TTL : 64      Window : 0      IP ID Class: ALL_ZERO
  Probability:  Linux 76%  ###############   macOS/BSD 24%  ####   Windows 0%
\end{logbox}

Under \texttt{--prefer-icmp}, the ICMP Echo Reply candidate independently confirms Linux at
67\% — a firewall-bypass path when closed-TCP RSTs are dropped but ICMP is answered.

%--------------------------------------------------------------------
\subsection{100-port scale test}\label{sub:scale}
\texttt{scan\_100\_ports.log} verifies the scanner at scale (two 50-port batches, honouring
\texttt{MAX\_PROBES}). Batch~1 finds the Windows RPC/SMB ports and fingerprints from the
SYN/ACK; Batch~2 (all closed) fingerprints from RST — both agree on Windows:
\begin{logbox}
   135 OPEN TTL=128 WIN=65535   139 OPEN TTL=128 WIN=8192
   445 OPEN TTL=128 WIN=65535   1025 OPEN TTL=128 WIN=65535

  Batch 1  Open-Port SYN/ACK  -> Windows 71%   (clock ~1001 Hz, uptime ~3d 5h)
  Batch 2  Closed-Port RST     -> Windows 69%   (IP ID STRICT_INCREMENTAL, WIN=0)
\end{logbox}

%--------------------------------------------------------------------
\subsection{Active defense on Windows Firewall}\label{sub:defwin}
\texttt{defense\_win\_test.log} exercises all four detection vectors on Windows, each
injecting a \texttt{netsh} block with a self-healing TTL. The fast-window vector trips on
the 8th SYN:
\begin{logbox}
[ALERT] 192.168.0.105 | fast port scan (SYN=8, distinct ports=8 within 10.0s)
[MITIGATION] Blocked TCP traffic from 192.168.0.105
             Firewall rule : CustomPortScanBlock_192_168_0_105 (Windows)
             Auto-release  : after 60s (self-healing)
\end{logbox}
The canary honeyport severs the attacker on the single packet that touches port~81, and the
invalid-flag vector catches the fingerprinter's own SYN+FIN probe instantly:
\begin{logbox}
[ALERT] 192.168.0.105 | hit canary honeyport (81)
[CRITICAL ALERT] 192.168.0.105 -> 192.168.0.116:21 |
        invalid TCP flag combination (SYN+FIN (Rule-Breaker), 0x03) on port 21
[MITIGATION] Released quarantine block on 192.168.0.105 (TTL 60s elapsed)
\end{logbox}
The slow/stealth vector separately caught a 2.5\,s-spaced scan at its 20th SYN. The injected
rule is the literal command \texttt{netsh advfirewall firewall add rule
name=CustomPortScanBlock\_192\_168\_0\_105 dir=in action=block remoteip=192.168.0.105
protocol=TCP profile=any}.

\paragraph{Live disruption (full-chain Run~2).}
With the defense active, the first exploit completed and every later stage timed out once
the rule landed — the attack chain was cut mid-execution:
\begin{logbox}
SYN 192.168.0.105 -> 192.168.0.116:55455  window_count=10  distinct_ports=5
[ALERT] Possible TCP port scan from 192.168.0.105
[MITIGATION] Blocked TCP traffic from 192.168.0.105
             Firewall rule: CustomPortScanBlock_192_168_0_105
  ... Exploit 2 -> Connection failed: timed out   traversal -> All 4 payloads failed ...
\end{logbox}

%--------------------------------------------------------------------
\subsection{Active defense on Linux Netfilter}\label{sub:defkali}
\texttt{defense\_kali\_test.log} mirrors every vector on Kali via
\texttt{iptables -I INPUT ... -j DROP}. The strongest Linux-specific evidence is the kernel
DROP counter climbing as the attacker's packets are dropped at the Netfilter layer:
\begin{logbox}
[ALERT] 192.168.56.1 | fast port scan (SYN=8, distinct ports=8 within 10.0s)
[MITIGATION] Firewall rule : CustomPortScanBlock_192_168_56_1 (Linux)
  iptables -L INPUT -n -v :
    12 pkts DROP tcp -- 192.168.56.1   ->   13 pkts   ->   14 pkts   ->   16 pkts
\end{logbox}

%--------------------------------------------------------------------
\subsection{SSH brute-force and post-exploitation}\label{sub:ssh}
Against a Windows OpenSSH host (\texttt{SSH-2.0-OpenSSH\_for\_Windows\_9.5}, fingerprinted
Windows 75\%), both SSH runs crack \texttt{test:test123} on attempt 15/20. Quiet mode stops
at recon; loud mode escalates to visible impact:
\begin{logbox}
  Status : SUCCESS   Logged in as test:test123 (attempt 15/20)
  whoami -> aisheepc\test   hostname -> AisheePC
  os_info -> Microsoft Windows [Version 10.0.26200.9445]
  [LOUD] desktop_proof  -> PROOF_OF_ACCESS.txt created on Desktop
  [LOUD] hacked_payloads-> HACKED.exe, HACKED.bat, hacked.py created on Desktop
  [LOUD] voice_alert    -> spooky voice alert played on victim speakers
\end{logbox}

%--------------------------------------------------------------------
\subsection{Mobile and BSD stacks}\label{sub:mobile}
The mobile targets confirm the engine generalises beyond desktop OSes. The iPad
(\texttt{ipad.log}) is attributed to macOS/BSD on its Darwin option order, and the Android
phone (\texttt{phn.log}) to Linux on its option layout and zero IP-ID:
\begin{logbox}
[iPad]    macOS / BSD 48%  options ['MSS','WSCALE','TIMESTAMP','SACK_PERMITTED'] (BSD layout)
[Android] Linux 58%        options ['MSS','SACK_PERMITTED','TIMESTAMP','WSCALE'] (Linux layout)
\end{logbox}

\begin{keypoint}[On honest reporting]
Two of the raw logs carry an over-optimistic ``99\%'' figure in a trailing prose summary,
but every actual fingerprint \emph{verdict block} reports 48--76\%. This report quotes only
the verdict blocks — the values the engine actually computes — never the inflated footer.
\end{keypoint}

\noindent The excerpts above are drawn from the twelve run logs, which are reproduced
\textbf{in full and verbatim in Appendix~C} — the complete, unedited proof of every claim in
this report.
